\section*{Lampiran A: Daftar Singkatan}
\phantomsection
\addcontentsline{toc}{section}{Lampiran A: Daftar Singkatan}

\begingroup
\setlength{\tabcolsep}{3pt}
\begin{longtable}{L{2.2cm} L{4.5cm} L{6.0cm}}
\caption{Daftar singkatan yang digunakan dalam laporan}
\label{tab:singkatan}\\
\toprule
\textbf{Singkatan} & \textbf{Kepanjangan} & \textbf{Arti atau Penggunaan}\\
\midrule
\endfirsthead
\toprule
\textbf{Singkatan} & \textbf{Kepanjangan} & \textbf{Arti atau Penggunaan}\\
\midrule
\endhead
AI      & Artificial Intelligence & Kecerdasan buatan yang digunakan untuk prediksi suhu secara otomatis.\\
API     & Application Programming Interface & Antarmuka pemrograman yang digunakan untuk membaca dan menulis data ke ThingSpeak.\\
DS18B20 & Dallas Semiconductor Temperature Sensor 18B20 & Sensor suhu digital berbasis protokol \textit{1-Wire} yang terhubung ke ESP32 pada pin GPIO4.\\
ESP32   & Espressif Systems ESP32 & Mikrokontroler dengan Wi-Fi \textit{built-in} yang digunakan sebagai perangkat utama pengumpul dan pengirim data suhu.\\
GPIO    & General Purpose Input/Output & Pin serbaguna pada ESP32; GPIO4 digunakan sebagai jalur komunikasi sensor DS18B20.\\
HTTP    & HyperText Transfer Protocol & Protokol komunikasi yang digunakan ESP32 dan Python untuk mengirim dan menerima data dari ThingSpeak.\\
IDE     & Integrated Development Environment & Lingkungan pengembangan perangkat lunak; Visual Studio Code digunakan untuk menulis skrip Python AI.\\
IoT     & Internet of Things & Paradigma konektivitas perangkat fisik ke internet; menjadi landasan arsitektur sistem pemantauan suhu ini.\\
JSON    & JavaScript Object Notation & Format data yang dikembalikan oleh API ThingSpeak saat Python mengambil riwayat suhu.\\
MAE     & Mean Absolute Error & Metrik evaluasi model yang mengukur rata-rata selisih absolut antara suhu aktual dan suhu prediksi.\\
ML      & Machine Learning & Cabang AI yang mempelajari pola data; digunakan untuk membangun model prediksi suhu berbasis Random Forest.\\
R\textsuperscript{2} & Coefficient of Determination & Metrik yang mengukur seberapa baik model menjelaskan variasi data suhu (R² Score).\\
RF      & Random Forest & Algoritma \textit{ensemble learning} berbasis banyak pohon keputusan yang digunakan sebagai model prediksi suhu.\\
RMSE    & Root Mean Squared Error & Metrik evaluasi model yang mengukur akar rata-rata kuadrat selisih antara suhu aktual dan prediksi.\\
SSID    & Service Set Identifier & Nama jaringan Wi-Fi yang digunakan oleh ESP32 dalam simulasi Wokwi, yaitu \texttt{Wokwi-GUEST}.\\
URL     & Uniform Resource Locator & Alamat web yang digunakan Python untuk mengakses endpoint API ThingSpeak.\\
VS~Code & Visual Studio Code & IDE yang digunakan untuk mengembangkan dan menjalankan skrip Python model AI Random Forest.\\
Wi-Fi   & Wireless Fidelity & Teknologi jaringan nirkabel yang digunakan ESP32 untuk terhubung ke internet dan mengirim data ke ThingSpeak.\\
\textit{1-Wire} & 1-Wire Communication Protocol & Protokol komunikasi serial satu kabel yang digunakan sensor DS18B20 untuk mengirim data suhu ke ESP32.\\
\bottomrule
\end{longtable}
\endgroup
